\selectlanguage{ngerman}%
\frenchspacing%
\begin{ledgroupsized}[r]{120mm}%
\footnotesize%
\pstart%
\noindent%
\textbf{Überlieferung:}%
\pend%
\end{ledgroupsized}%
\begin{ledgroupsized}[r]{114mm}%
\footnotesize%
\pstart%
\parindent -6mm%
\makebox[6mm][l]{\textit{L}}%
Aufzeichnung:
LH XXXVIII Bl.\;182.
Ein Zettel (6\;x\;10 cm),
Ränder unregelmäßig beschnitten und zum Teil abgerissen.
Zwei vollbeschriebene Seiten.
\pend%
\end{ledgroupsized}%
\vspace{5mm}%
\begin{ledgroup}%
\footnotesize%
\pstart%
\noindent%
\textbf{Datierungsgründe:}
Die Datierung der Aufzeichnung N.\;\ref{55813-C} beruht auf %
inhaltlichen Betrachtungen.
Siehe die ausführliche Begründung auf S.\;\pageref{55813-C_dat}.%
\pend%
\end{ledgroup}%
\selectlanguage{latin}%
\frenchspacing%
\vspace{8mm}%
\normalsize%
\pstart%
\noindent%
\lbrack%
\href{\Scan{LH_38_182r}}%
{182\,r\textsuperscript{o}}%
\rbrack\
Consideratio incidit magni in Mechanicis Momenti,%
\indexS{consideratio magni momenti in mechanicis}%
\indexS{mechanica}
cur caeteris paribus melius
\edtext{%
sit idem punctum moventis%
}{%
\lemma{sit}%
\Bfootnote{%
\textit{(1)}~movens %
\textit{(2)}~idem punctum moventis%
~\textit{L}%
}}
eidem puncto movendi insistere,
quam ab alio ad aliud procedi.
\pend%
\begin{edfigure}[2.05em][1.5em]
\includegraphics%
{gesamttex/edit_VIII-4/images/LH_38_182_d1_182r.pdf}
\edcaption[0.25em]{%
\lbrack\textit{Fig.\:1}\rbrack%
}%
\end{edfigure}%
\pstart%
\noindent%
Exempli causa
si centro \textit{A}
radius \textit{AC} %
impellere debeat %
radium \textit{De}%
\indexS{radius}
\edtext{%
centri \textit{D},
is corporis movendi punctum \textit{e} %
primum impellet%
}{%
\lemma{centri \textit{D},}%
\Bfootnote{%
\textit{(1)}~impellet %
\textit{(2)}~is %
\textit{(a)}~punctum movendi %
\textit{(b)}~corporis movendi %
\lbrack...\rbrack\ primum impellet%
~\textit{L}%
}}
puncto suo \textit{B},
deinde puncto suo \textit{C};
ubi,
ut frictionis%
\indexS{frictio}
mentionem non faciam
(quam patitur \textit{e}
incedendo a \textit{B} ad \textit{C},
quae ex parte saltem
\edtext{%
rotula%
\indexS{rotula}
vel trochlea%
\indexS{trochlea}
tollitur),%
}{%
\lemma{rotula}%
\Bfootnote{%
\textit{(1)}~tollitur, %
\textit{(2)}~vel trochlea tollitur),%
~\textit{L}%
}}
sciendum est,
vi impulsus%
\indexS{vis impulsus}%
\indexS{impulsus}
nonnihil corporum partes cedere et comprimi;%
\indexS{compressio}
itaque primum partes radii%
\indexS{radius}
sive brachii \textit{AC}%
\indexS{brachium}
ad punctum \textit{B} comprimentur,
et digresso resistente
\edtext{%
\textit{e} a \textit{B} ad \textit{C},%
}{%
\lemma{\textit{e} a}%
\Bfootnote{%
\textit{(1)}~\textit{C} ad %
\textit{(2)}~\textit{B} ad \textit{C},%
~\textit{L}%
}}
iterum se restituent;%
\indexS{restitutio}
dum interea partes comprimentur
ad punctum \textit{C}.
Omnis autem
\edtext{%
vis quam
\lbrack%
\href{\Scan{LH_38_182v}}%
{182\,v\textsuperscript{o}}%
\rbrack\
ipsum corpus quodammodo sentit,%
\indexS{vis sentita a corpore}
et quae in%
}{%
\lemma{vis}%
\Bfootnote{%
\textit{(1)}~in \lbrack\href{\Scan{LH_38_182v}}{182~v\textsuperscript{o}}\rbrack\ %
\textit{(2)}~quam ipsum %
\lbrack...\rbrack\ quae in%
~\textit{L}%
}}
partes corporum
(licet insensibiliter)
\ledpb%
comprimendas impensa est,%
\indexS{vis impensa in partes comprimendas}
periit.%
\indexS{vis perita}
At
\edtext{%
si punctis semper iisdem%
}{%
\lemma{si}%
\Bfootnote{%
\textit{(1)}~eidem semper puncto %
\textit{(2)}~punctis semper iisdem%
~\textit{L}%
}}
tota vis%
\indexS{vis tota}
innitatur%
\lbrack,\rbrack\
non nisi
\edlabel{LH_38_182_zfk-1}%
\edtext{}{%
\xxref{LH_38_182_zfk-1}{LH_38_182_zfk-2}%
\lemma{semel}%
\Bfootnote{%
\textit{(1)}~toto tractu %
\textit{(2)}~fiet compressio %
\lbrack...\rbrack\ continuetur pressio, %
\textit{(a)}~non tamen perinde con %
\textit{(b)}~toto tractu %
\lbrack...\rbrack\  sed mox cessabit.
\textit{(aa)}~Caeterum putem tractu temporis illam \textlangle dis\textrangle\ %
\textit{(bb)}~Et hoc %
\lbrack...\rbrack\ et restitutionem%
~\textit{L}%
}}%
semel fiet compressio%
\indexS{compressio}
et restitutio;%
\indexS{restitutio}
adeoque minus virium perdetur.%
\indexS{vis perdita}
Licet enim continuetur pressio,%
\indexS{pressio}
toto tractu%
\indexS{tractus}
vel impulsu%
\indexS{impulsus}
durante,
non tamen continuabitur compressio,%
\indexS{compressio}
sed mox cessabit.
\edlabel{LH_38_182_182v_fodinae-1}%
Et hoc contingit,
v.\:g.
cum ansa recurva parallelogramma%
\indexS{ansa recurva parallelogramma}
apud fodinas moveatur,%
\indexS{fodina}
ubi puncta aggressionis%
\indexS{punctum aggressionis}
\textit{a} et \textit{b}
sunt semper eadem.%
\edlabel{LH_38_182_182v_fodinae-2}%
\pend%
\begin{edfigure}[1.90em][1.55em]
\includegraphics%
{gesamttex/edit_VIII-4/images/LH_38_182_d2_182v.pdf}
\edcaption[-0.40em]{%
\lbrack\textit{Fig.\:2}\rbrack%
}%
\end{edfigure}%
\pstart%
Puto tamen tractu temporis%
\indexS{tractus temporis}
minui compressionem illam%
\indexS{compressio}
et restitutionem%
\indexS{restitutio}%
\edlabel{LH_38_182_zfk-2}
in casu priori,
quia semper eodem tendit nisus,%
\indexS{nisus}
corpus autem saepe eodem modo compressum%
\indexS{corpus compressum}%
\indexS{compressio}
non amplius facile restituitur%
\indexS{corpus restitutum}%
\indexS{restitutio}
(et hanc inter caeteras causam puto,%
\indexS{causa inter caeteras}
cur initio novae machinae magis resistant).%
\indexS{machina nova magis resistens}
Posteriori vero casu
tardius opinor desinit,
quia in
\edtext{%
tantum trahuntur iterum%
}{%
\lemma{tantum}%
\Bfootnote{%
\textit{(1)}~retrahitur iterum %
\textit{(2)}~trahuntur iterum%
~\textit{L}%
}}
partes corporum%
\indexS{pars corporis tracta}
in quantum erant impulsae%
\indexS{pars corporis impulsa}
et compressae.%
\indexS{pars corporis compressa}
Et quoniam rotae dentatae%
\indexS{rota dentata}
ad priorem casum pertinent,
quae tamen satis utiles deprehenduntur,
forte minus est malum
quam prima fronte videtur.
\pend%
\newpage%
\pstart%
\noindent%
\centering%
*     *     *
\pend%
\vspace{3.0em}%
\selectlanguage{german}%
\frenchspacing%
\footnotesize%
\pstart%
\noindent%
\textls{Ausführliche Datierungsgründe:}%
\pagelabel{55813-C_dat}
Die titellose und undatierte Aufzeichnung N.\;\ref{55813-C}, in deren Mittelpunkt mögliche praktische Folgen einer allgemeinen Betrachtung zur Kraftanwendung stehen, weist weder Spuren eines Wasserzeichens noch besondere äußere Merkmale auf, die eine klare chronologische Einordnung ermöglichen würden.
Auch der Inhalt bietet kaum Anhaltspunkte, um die Entstehungszeit der Aufzeichnung zu bestimmen.
Eine Beobachtung im Text (\refpassage{LH_38_182_182v_fodinae-1}{LH_38_182_182v_fodinae-2}) zeigt jedoch, dass Leibniz zu dem Zeitpunkt, als er sie verfasst hat, bereits eine gewisse Erfahrung mit technischen Vorrichtungen im Bereich des Bergbaus haben musste.
Daraus lässt sich schließen, dass die Aufzeichnung N.\;\ref{55813-C} vermutlich nicht entstanden ist, bevor Leibniz im Herbst 1680 begonnen hatte, seine Projekte für die Verbesserung des Bergbaus im Harz\indexO{Harz} regelmäßig zu betreuen.
Eine frühere Datierung ist trotzdem möglich.
Das Papierformat deutet jedenfalls auf eine Entstehung nach Leibnizens Umzug von Paris\indexO{Paris} nach Hannover\indexO{Hannover} gegen Ende 1676.
\pend%
\selectlanguage{latin}%
\frenchspacing%
\normalsize%
\newpage%
